\section{Training Budget and Recovery at Larger Dimensions}
\label{sec:frontier}

\textbf{Longer training.} Several configurations that failed at their
initial training budget succeeded with more steps. For composition
at $K{=}8$, one of five seeds converged at 20K steps and four of five
at 40K.
% <!-- evidence: R8 -->
At $K \in \{12, 16\}$, none of five seeds converged at 40K; at 80K,
three of three and two of three fresh seeds converged, respectively.
% <!-- evidence: R6 -->
In the dimension sweep, configurations at $d \geq 32$ failed at an
8K-step budget but began learning at 6--16K steps when training was
extended to 20K.
% <!-- evidence: R8 -->
These reversals motivate re-testing unsuccessful configurations at
2--2.5$\times$ the initial budget before drawing conclusions about
trainability.
% <!-- evidence: R8 -->
An extension does not guarantee recovery: the remaining unsuccessful
$K{=}16$ seed still failed at 120K steps,
% <!-- evidence: R6 -->
and the $d{=}32$ configurations below plateaued below the
pre-registered success threshold.
% <!-- evidence: R8 -->

\Needspace{9\baselineskip}
\textbf{Recovery at larger dimensions.} With $K = d/4$ and encoder
width fixed at 64, mean recovery cosine declines as $d$ increases.
All reported runs had plateaued: at $d{=}32$, three seeds reached
0.877/0.909/0.915 after 100K steps; at $d{=}48$, one seed reached
0.7196 after 100K; at $d{=}64$, one seed reached 0.3882 after 150K.
Their $\recninety$ scores were 0.413/0.632/0.653 at $d{=}32$
(all below the pre-registered 0.7 success threshold), 0.002 at
$d{=}48$, and 0.0 at $d{=}64$.
% <!-- evidence: R7 -->
At $d{=}32$, learned effective rank reaches 91--97\% of the target,
but recovery remains below the success threshold. High effective rank
alone does not ensure accurate bindings.
% <!-- evidence: R7 -->
